% AUTHORED
\section{Multi-head attention}\label{sec:10}
\conceptbadge{10}

A single attention head computes one set of query-key matches and one weighted
sum of values. Multi-head attention repeats that operation several times with
different learned projections. The outputs are concatenated and mixed with a
final linear projection.

\subsection*{Why multiple heads?}
One head has one similarity space. Multiple heads allow the model to attend to
different kinds of information at the same time: nearby syntax, repeated names,
delimiter structure, or longer-range dependencies. The heads are not told what
to specialize in; specialization emerges because each head has separate
projection matrices and receives gradients through the loss.

\subsection*{Shapes}
Let the model width be $C$ (the channel dimension $C$ from \S00) and the
number of heads be $h$. A common choice is the per-head key/query width
$d_k = C/h$. For each head $r$, with input $X \in \mathbb{R}^{B \times T \times C}$
from the previous block,
\[
  Q_r = XW^Q_r,\qquad K_r = XW^K_r,\qquad V_r = XW^V_r,
\]
then
\[
  O_r = \operatorname{softmax}\left(\frac{Q_rK_r^\top}{\sqrt{d_k}}\right)V_r.
\]
The head outputs $O_1,\dots,O_h$ are concatenated along the channel dimension,
returning to width $C$, and an output projection mixes information across heads:
\[
  O = \operatorname{concat}(O_1,\dots,O_h)W_O.
\]

\subsection*{Cost and capacity}
Using $h$ heads does not necessarily multiply the channel width, because each
head usually uses a smaller internal dimension. The total attention score cost
still scales with sequence length squared. More heads increase representational
flexibility, but too many tiny heads can become inefficient or underpowered.

\subsection*{Decoder setting}
In a decoder-only GPT, every head is causal. The same lower-triangular mask from
concept 09 applies inside each head before softmax. P6 builds the multi-head
wrapper around the attention function you already implemented, and P8 places it
inside a full transformer block. Multi-head attention still assumes that each
token vector already carries information about where it appears in the sequence;
the next module adds that missing position signal.
